\chapter[Experimentos complementares: Síntese de Idade e ConvNeXt-ViT]{Experimentos complementares: Síntese de Idade e ConvNeXt-ViT}
\label{ap:modelos-perifericos}

Este apêndice registra dois experimentos implementados ao longo do projeto, mas que não compõem o eixo central de comparação do Capítulo~\ref{cap:Metodologia} e do Capítulo~\ref{cap:resultados}. A função do apêndice é documentar decisões de escopo: o Modelo com Síntese de Idade fica como referência histórica herdada da versão de 2018 do pré-projeto, enquanto o ConvNeXt-ViT fica como estudo de capacidade de uma fusão híbrida ConvNeXt mais ViT. Assim, os resultados principais permanecem concentrados nos modelos avaliados de forma comparável em FIW.

\section{Modelo com Síntese de Idade}

O Modelo com Síntese de Idade é a linha de base histórica do projeto. Ele usa uma ResNet-18 com pesos do ArcFace, pré-treinada em reconhecimento facial, em uma configuração siamesa para comparação de pares. A ideia original era usar o módulo de síntese de idade SAM \cite{alaluf2022sam} para gerar variantes etárias das faces e reduzir a diferença aparente de idade entre parentes. Na prática, o SAM exige pesos pré-treinados específicos de cerca de 4,5 GB, e parte das execuções históricas ocorreu sem esse módulo ativo. Por isso, esse modelo é tratado aqui apenas como registro de continuidade do projeto, não como competidor da comparação principal.

A motivação do experimento é plausível: em relações entre pais e filhos, ou entre avós e netos, duas pessoas biologicamente próximas podem aparecer em fases muito diferentes da vida. Quando disponível, o protocolo com SAM geraria versões criança, jovem e idosa de cada face e combinaria o escore do par original com escores auxiliares dessas variantes. Como essa manipulação altera o dado de entrada e não foi usada de forma uniforme nas execuções finais, ela foi mantida fora do eixo principal do trabalho.

O resultado registrado para esse modelo foi obtido em KinFaceW-I, com AUC 0,682 e acurácia de 68,5\%. Ele é mantido apenas como referência de partida, pois não competiu no benchmark principal em FIW. O ViT-FaCoR, o DINOv2-LoRA, o Modelo com Recuperação, o AdaFace-Cosseno e o AdaFace-Regional operam sobre as faces originais alinhadas, sem gerar variantes etárias.

\section{ConvNeXt-ViT, híbrido ConvNeXt e ViT}

O ConvNeXt-ViT combina dois backbones distintos, ConvNeXt-B \cite{liu2022convnet} e vit\_base\_patch16\_224, sob a hipótese de que atributos locais convolucionais e atributos globais de Transformer são complementares no domínio facial. Os dois backbones operam em paralelo sobre a mesma imagem e produzem representações que são fundidas em um módulo aprendível antes da projeção final. A fusão suporta quatro modos. O padrão é a atenção, com alternativas em concatenação, gating multiplicativo e fusão bilinear. A saída final é um embedding de 512 dimensões com normalização $L^2$. O modelo total tem cerca de 176 milhões de parâmetros, mais que o dobro do ViT-FaCoR, e essa expansão de capacidade é a motivação central de mantê-lo como estudo complementar.

O ConvNeXt-ViT atinge AUC 0,850 em FIW, valor idêntico ao do ViT-FaCoR com 86 milhões de parâmetros. A leitura central desse resultado é que dobrar a capacidade do backbone não traz ganho proporcional em AUC neste banco, o que confere ao ConvNeXt-ViT papel argumentativo complementar sem necessidade de ocupar lugar no eixo central de comparação.

\section{Estudo de congelamento do ConvNeXt-ViT}

O ConvNeXt-ViT, por ter 176 milhões de parâmetros e dois backbones grandes, oferece um caso natural para investigar como o congelamento de pesos pré-treinados afeta o desempenho final. As quatro estratégias consideradas são tratadas como variantes do mesmo modelo. A variante mais permissiva treina todos os pesos desde o início, sem congelamento. No outro extremo, os dois backbones permanecem congelados durante todo o treino e a otimização atinge só a fusão e a cabeça, o que reduz bastante o custo computacional por época. Entre os dois extremos, uma variante começa com backbones congelados e libera todos os pesos depois de algumas épocas, treinando até saturar a métrica de validação. A última descongela apenas as camadas superiores de cada backbone, em geral os dois últimos estágios do ConvNeXt e os quatro últimos blocos do ViT, com taxa de aprendizado reduzida por um fator entre 1e-2 e 1e-3 em relação à cabeça.

As execuções consolidadas do ConvNeXt-ViT cobrem essas estratégias em diferentes combinações. A Tabela~\ref{tab:convnext_vit_freeze_appendix} resume os resultados usados apenas como apoio interpretativo. O maior AUC do modelo (0,850) ocorreu sem congelamento, indicando para esta arquitetura que aumentar a capacidade do backbone não trouxe ganho sobre o ViT-FaCoR, apesar de quase dobrar o número de parâmetros.

\begin{table}[h]
\centering
\small
\caption{Resumo do estudo de congelamento do ConvNeXt-ViT.}
\label{tab:convnext_vit_freeze_appendix}
\begin{tabular}{lll}
\hline
\textbf{Estratégia} & \textbf{AUC teste} & \textbf{Interpretação} \\
\hline
sem congelamento & 0{,}850 & melhor resultado da família \\
sem congelamento (repetição) & 0{,}845 & confirma teto próximo ao melhor resultado \\
backbones congelados & 0{,}823 & cabeça sozinha insuficiente \\
congelar e descongelar tudo & 0{,}813 & execução interrompida por OOM \\
congelar e descongelar tudo, lote menor & 0{,}848 & quase empata, mas com sobreajuste após descongelamento \\
descongelamento parcial & 0{,}828 & taxa do backbone conservadora demais \\
\hline
\end{tabular}
\end{table}
